\documentclass[11pt]{article}

\usepackage[margin=1in]{geometry}
\usepackage{amsmath,amssymb}
\usepackage{booktabs}
\usepackage{enumitem}
\usepackage{hyperref}
\hypersetup{hidelinks}

\title{Scaling Book Exercises: Section 7}
\author{Lucas Yan}
\date{}

\begin{document}
\maketitle

\noindent
\textbf{Chapter:} All About Transformer Inference\\
\textbf{Source scan:} \texttt{../handwritten/section\_07\_handwritten.pdf}\\
\textbf{Reference:} \url{https://jax-ml.github.io/scaling-book/inference/}

\paragraph{Model.}
The worked problems use
\[
L=64,\quad D=4096,\quad F=16384,\quad N=32,\quad
K=8,\quad H=256,\quad V=32128.
\]
Here $L$ is the number of layers, $D$ is the model dimension, $F$ is the
feed-forward dimension, $N$ is the number of query heads, $K$ is the number
of KV heads, $H$ is the head dimension, and $V$ is the vocabulary size.

\section*{Exercise 1}

The gated MLP has three $D\times F$ matrices per layer, so
\[
P_{\mathrm{MLP}}=3LDF
=3(64)(4096)(16384)
=12{,}884{,}901{,}888.
\]
The attention matrices contribute
\begin{align*}
P_{\mathrm{attention}}
&=L\bigl(DNH+DKH+DKH+DNH\bigr)\\
&=2LDH(N+K)\\
&=2(64)(4096)(256)(32+8)\\
&=5{,}368{,}709{,}120.
\end{align*}
Because the input and output embedding matrices are shared,
\[
P_{\mathrm{vocab}}=DV=(4096)(32128)=131{,}596{,}288.
\]
Therefore,
\begin{align*}
P
&=3LDF+2LDH(N+K)+DV\\
&=18{,}385{,}207{,}296\\
&\approx\boxed{18.4\text{B parameters}}.
\end{align*}

For each token, every layer caches one key and one value for each KV head.
Thus the KV cache contains
\[
2LKH=2(64)(8)(256)=262{,}144
\]
elements per token. With int8 elements,
\[
\boxed{262{,}144\text{ bytes/token}
\approx262\text{ kB/token}=256\text{ KiB/token}}.
\]

\section*{Exercise 2}

A TPU v5e $4\times4$ slice contains 16 chips, each with 16 GB of HBM:
\[
M_{\mathrm{slice}}=16(16\text{ GB})=256\text{ GB}.
\]
The int8 model weights occupy approximately 18.4 GB. A 128,000-token KV
cache for one sequence occupies
\[
(262{,}144\text{ bytes/token})(128{,}000\text{ tokens})
=33.55\text{ GB}.
\]
The maximum batch size is therefore
\[
B_{\max}
=\left\lfloor
\frac{256-18.4}{33.55}
\right\rfloor
=\boxed{7}.
\]
If $K$ falls from 8 to 1, each KV cache is eight times smaller, so
\[
\boxed{B_{\max}\approx 8(7)=56}.
\]

\section*{Exercise 3}

The ideal parameter-loading time is bytes divided by aggregate HBM bandwidth.
Each TPU v5e supplies approximately $8.2\times10^{11}$ bytes/s, so the
$4\times4$ slice supplies
\[
16(8.2\times10^{11})=1.312\times10^{13}\text{ bytes/s}.
\]
For 18.4 GB of int8 parameters,
\[
T_{\mathrm{parameters}}
=\frac{18.4\times10^9}{1.312\times10^{13}}
\approx1.4\times10^{-3}\text{ s}
=\boxed{1.4\text{ ms}}.
\]
This is an optimistic lower bound assuming perfect sharding and full HBM
bandwidth utilization.

\section*{Exercise 4}

\subsection*{Sharding}

A TPU v5e $4\times4$ slice has two physical ICI dimensions, each four chips
wide. For prefill, the approximate tensor-parallel roofline is
\[
p_{\mathrm{prefill}}\approx\frac{F}{2200}
=\frac{16384}{2200}\approx7.4.
\]
A natural mapping therefore uses one four-chip axis for tensor parallelism
and the other for sequence parallelism.

Decode is bandwidth-bound, so tensor parallelism can extend further. Let
\[
\beta=\frac{W_{\mathrm{HBM}}}{W_{\mathrm{ICI}}}\approx8.
\]
Comparing parameter-loading time with ICI communication gives
\[
p_{\mathrm{decode}}\lesssim\frac{F}{B\beta}.
\]
For $B=7$,
\[
\frac{F}{B\beta}
=\frac{16384}{7(8)}
\approx293,
\]
which is much larger than the available 16 chips. Thus a 16-way
model-parallel decode is below this idealized bandwidth roofline.

The KV cache has shape
\[
\mathrm{KV}[2,B,S,K,H].
\]
Shard KV heads first and batch second:
\[
\mathrm{KV}[2,B_Z,S,K_Y,H].
\]
One natural $4\times4$ layout uses $Y=4$ and $Z=4$, distributing the KV
cache across all 16 chips. This layout requires two AllToAll operations per
attention layer to move query and output activations between model-sharded
and batch-sharded layouts.

\subsection*{Generation latency}

For $B=7$ and a 128,000-token context, the KV-cache loading time is
\[
T_{\mathrm{KV}}
=\frac{7(33.55\times10^9)}
{16(8.2\times10^{11})}
\approx17.9\text{ ms}.
\]
The parameter-loading time is 1.4 ms. Using approximately
$3.93\times10^{14}$ int8 operations/s per chip, the compute time is
\[
T_{\mathrm{compute}}
=\frac{2BP}{16C}
=\frac{2(7)(18.4\times10^9)}
{16(3.93\times10^{14})}
\approx0.041\text{ ms}.
\]
Parameter loading dominates the linear layers, so the theoretical step time is
\begin{align*}
T_{\mathrm{step}}
&=T_{\mathrm{KV}}
  +\max(T_{\mathrm{parameters}},T_{\mathrm{compute}})\\
&\approx17.9+1.4\\
&=\boxed{19.3\text{ ms}}.
\end{align*}
This excludes fixed collective latency and other runtime overheads.

\section*{Exercise 5}

Suppose the MLP becomes a mixture of experts with $E=16$ total experts and
$k=2$ activated experts per token.

\subsection*{Parameters}

The total parameter count is
\begin{align*}
P_{\mathrm{total}}
&=3ELDF+2LDH(N+K)+DV\\
&=211{,}658{,}735{,}616\\
&\approx\boxed{212\text{B}}.
\end{align*}
Only $k=2$ experts are used by a given token, so the activated parameter count
is
\begin{align*}
P_{\mathrm{activated}}
&=3kLDF+2LDH(N+K)+DV\\
&=31{,}270{,}109{,}184\\
&\approx\boxed{31.2\text{B}}.
\end{align*}

\subsection*{Critical batch size}

The MoE stores $E$ times as many MLP weights but performs only $k$ times as
much MLP computation per token. Therefore its critical batch size increases
by approximately $E/k$. Starting from the bf16 TPU v5e dense-model estimate
of 240 tokens,
\[
B_{\mathrm{crit}}
\approx240\frac{E}{k}
=240\frac{16}{2}
=\boxed{1920\text{ tokens}}.
\]

\subsection*{KV cache and forward FLOPs}

MoE changes the MLP, not attention, so the KV cache remains
\[
\boxed{2LKH=262{,}144\text{ elements/token}},
\]
or approximately 262 kB/token in int8.

For $T$ tokens, the forward-pass FLOP count is approximately
\[
2P_{\mathrm{activated}}T
\approx2(31.2\times10^9)T
=\boxed{62.4\times10^9T\text{ FLOPs}}.
\]

\end{document}
